Sanguosuo follows a three-tier client--server--database architecture. The Android client communicates with a REST API backend over HTTP; the backend reads and writes to a relational MySQL database. During local development, Docker Compose orchestrates the backend and database as isolated containers on the developer's machine.

\subsection{High-Level Architecture}

\begin{center}
\adjustbox{max width=\linewidth}{%
\begin{tabular}{llp{6cm}}
\toprule
\textbf{Tier} & \textbf{Technology} & \textbf{Role} \\
\midrule
Client (Android) & Kotlin + Jetpack Compose & UI, state management, API calls \\
Backend (API)    & Node.js + Express (TypeScript) & Business logic, authentication, data access \\
Database         & MySQL 8                        & Persistent storage of all game and user data \\
Infrastructure   & Docker Compose                 & Local orchestration of API and DB containers \\
\bottomrule
\end{tabular}%
}
\end{center}

\subsection{Android Client}

The client is structured around the MVVM (Model--View--ViewModel) pattern encouraged by Jetpack:

\begin{itemize}
  \item \textbf{UI layer:} Composable screens (Welcome, Sign-Up pages 1 and 2, Sign-In, Home, Hero detail, Profile, Search) written entirely in Jetpack Compose. Navigation between screens is managed by a single \texttt{NavHost} defined in \texttt{ui/navigation/}.
  \item \textbf{ViewModel layer:} Each screen has a corresponding \texttt{ViewModel} that holds \texttt{StateFlow} objects. The UI collects these flows and re-composes on state changes. Business decisions such as input validation are made here.
  \item \textbf{Repository layer:} Repositories under \texttt{data/repository/} are the single source of truth for data. They call Retrofit service interfaces to fetch remote data and expose results to ViewModels.
  \item \textbf{Remote layer:} Retrofit interfaces and their corresponding data-transfer objects (DTOs) under \texttt{data/remote/} map HTTP responses to Kotlin data classes using Kotlinx Serialization.
  \item \textbf{Session layer:} The \texttt{data/session/} module uses Jetpack DataStore Preferences to persist the authentication token across app restarts.
  \item \textbf{DI:} A lightweight dependency-injection module under \texttt{di/} manually provides shared instances (Retrofit client, repositories) to avoid tight coupling.
\end{itemize}

\subsection{Backend API}

The backend is a Node.js Express application written entirely in TypeScript. Its internal layout mirrors a classic layered web-server design:

\begin{itemize}
  \item \textbf{Routes} (\texttt{routes/}) declare URL paths and attach middleware chains.
  \item \textbf{Middleware} (\texttt{middleware/}) performs cross-cutting concerns: JSON content-type checking, HTTP method validation (\texttt{ValidatorMiddleware}), centralized error handling (\texttt{errorHandler}), request logging (Morgan), and fixed-path callbacks for the root and health endpoints.
  \item \textbf{Controllers} (\texttt{controllers/}) parse and validate request parameters before delegating to the service layer.
  \item \textbf{Services} (\texttt{services/}) contain business logic: user registration, login, and hero lookup. Services call the database connection pool.
  \item \textbf{Connections} (\texttt{connections/DbConnection.ts}) wraps \texttt{mysql2}'s \texttt{createPool} call, configured for a maximum of five concurrent connections.
  \item \textbf{Helpers and types} provide shared response-formatting utilities and TypeScript type definitions.
\end{itemize}

The API exposes the following endpoints:

\begin{center}
\adjustbox{max width=\linewidth}{%
\begin{tabular}{llp{7cm}}
\toprule
\textbf{Method} & \textbf{Path} & \textbf{Description} \\
\midrule
\texttt{POST} & \texttt{/api/auth/register} & Create a new user account \\
\texttt{POST} & \texttt{/api/auth/login}    & Authenticate and receive a session token \\
\texttt{GET}  & \texttt{/api/heroes/id/:heroId}     & Retrieve a hero by its unique identifier \\
\texttt{GET}  & \texttt{/api/heroes/name/:heroName} & Retrieve a hero by name (partial match) \\
\texttt{GET}  & \texttt{/health}            & Health check (returns server and DB status) \\
\texttt{GET}  & \texttt{/}                  & API root information \\
\bottomrule
\end{tabular}%
}
\end{center}

\subsection{Database}

The MySQL database is organized into three conceptual groups of tables:

\begin{itemize}
  \item \textbf{Lookup / enumeration tables:} \texttt{UserStatus}, \texttt{UserRole}, \texttt{HeroFaction}, \texttt{HeroGender}, \texttt{HeroRole}, \texttt{SkillTag}, \texttt{CardType}, \texttt{CardSuit}. These store coded values and human-readable names.
  \item \textbf{Core entity tables:} \texttt{User}, \texttt{Hero}, \texttt{HeroSkill}, \texttt{Card} (and sub-type tables \texttt{BasicCard}, \texttt{ToolCard}, \texttt{WeaponCard}, \texttt{ArmorCard}, \texttt{HorseCard}, \texttt{TreasureCard}---these will be implemented in a future version), \texttt{Session}, \texttt{RevisionText}, \texttt{Revision}, \texttt{Article}, \texttt{Image}.
  \item \textbf{Relationship / bridge tables:} \texttt{HeroSkillTag}, \texttt{HeroCombo}, \texttt{HeroBelongsFaction}, \texttt{HeroBelongsRole}, \texttt{CardSuitAndRank}, \texttt{HeroSaves}.
\end{itemize}

The schema enforces data integrity through primary keys, foreign keys, unique constraints, and \texttt{CHECK} constraints, including a regular-expression constraint that rejects any stored password that does not match the bcrypt hash format, and a SHA-224 content-hash constraint on \texttt{RevisionText}.

\subsection{Request Flow: Registration Example}

\begin{enumerate}
  \item The Android client validates email format and phone-number shape locally, confirms the two password fields match, and sends a JSON \texttt{POST} to \texttt{/api/auth/register}.
  \item The \texttt{ValidatorMiddleware} checks the HTTP method and \texttt{Content-Type} header.
  \item \texttt{AuthController} extracts and checks required fields.
  \item \texttt{AuthService} generates an eight-digit user ID, looks up the default status and role, hashes the password with \texttt{bcrypt} at cost 10, and inserts the account via a parameterized query through the connection pool.
  \item The controller returns \texttt{201 Created} with a success envelope.
\end{enumerate}

Hero lookup follows a similar path through \texttt{HeroController} and \texttt{HeroService}, where the service joins the \texttt{Hero}, \texttt{HeroFaction}, \texttt{HeroSkill}, and \texttt{HeroSkillTag} tables using parameterized queries and assembles the JSON response.

\subsection{Infrastructure}

Docker Compose defines two services: \texttt{sanguosuo-server} (the Express API) and \texttt{sanguosuo-db} (MySQL). The API container is built from the \texttt{server/Dockerfile}, mounts the source directory as a volume so that Nodemon hot-reloads on file changes during development, and exposes port 8080. The database container uses the official \texttt{mysql:oraclelinux9} image, persists data in a named volume, and mounts the \texttt{db/} directory at \texttt{/sql} read-only so that schema and seed files are accessible inside the container.
